% Pursuing Stacks (cotes 134-2 to 134-8): archive cote 134-2, site page 48 (PDF page 49).
% appendix (letters to Breen). Draft compared: ps1.tex:1226-1248.
% AI-transcribed by Claude Opus 5.5 (claude-opus-5-5, adaptive thinking, effort high), Message Batches API, batch ps1 (msgbatch_01Rv5hCaNTtUt4JHx1yqoHPL), 2026-11-12; not yet checked by eye.
\scanpage{134-2/48}
\mspage{33}

\deleted{\uline{Transcription d'une lettre manuscrite de A. Grothendieck}}

Ville\hand{cun} le 17.2.1975

Cher Larry,

\margin{(\unsure{App.6|Appr.6})} Encore un ``afterthought'' à une lettre-fleuve  sur le yoga homotopique.
Comme tu sais sans doute, à un topos X on associe canoniquement un \xout{xxxxxxxxxx}
pro-ensemble simplicial, donc un``pro-type  d'homotopie'' en un sens convenable.
Dans le cas où X est ``localement homotopiquement trivial'', le pro-objet
associé est essentiellement constant en tant que pro-objet dans la catégorie
homotopique, donc X définit un objet de la catégorie homotopique usuelle,
qui est son\hand{|}``type d'homotopie''.  De même, si X est  ``localement homotopique-%
ment  trivial en dim \hand{$\leq$} n '', il définit un type d'homotopie  ordinaire ``tron-%
qué en dim \hand{$\leq$} n''  -  construction familière pour i = 0 ou 1, même à des
gens comme moi qui ne connaissent guère l'homotopie !

Ces constructions sont fonctorielles en X. D'ailleurs, si  f: X $\longrightarrow$ Y
est un morphisme de topos,  Artin-Mazur ont donné  une condition nécéssaire
et suffisante  \uline{cohomologique}   pour que ce soit  une\hand{|}``équivalence d'homo-%
topie en dim \hand{$\leq$} n\hand{\unsure{''}}: c'est que  $H^{i}(Y,F) \xrightarrow{\hand{$\sim$}} H^{i}(X,f^{*}(F))$   pour i \hand{$\leq$} n, et
tout faisceau de groupes  \uline{localement constant}  F sur Y, en se restreignant
de plus à  i\hand{$\leq$} 1  dans le cas \unsure{F} non commutatif. Ce critère, en termes
de n-gerbes ``localement constantes'' \hand{$\mathcal{F}$}  sur Y, s'interprète par la condition
que  \hand{$\mathcal{F}$}(Y) $\longrightarrow$\hand{$\mathcal{F}$}(X)   est une n-équivalence  pour tout tel \hand{$\mathcal{F}$} et  i\hand{$\leq$} n.
Il est certainement vrai que ceci équivaut encore au critère suivant

(A)  Pour tout n-champ  ``localement constant'' \hand{$\mathcal{F}$}  sur Y, le
n-foncteur  \hand{$\mathcal{F}$}(Y) $\longrightarrow$ f$^{*}$(\hand{$\mathcal{F}$})(X)  est une n-équivalence

ou encore à

(B)  Le n- foncteur \hand{$\mathcal{F}$}$\longrightarrow$ f$^{*}$(\hand{$\mathcal{F}$})   allant de la n-catégorie
des (n-1) - champs localement constants  sur Y dans celle des (n-1) -champs
localement  constants sur X,  est une n-équivalence.

En d'autres termes, les constructions sur un topos X  qu'on peut
faire en termes de (n-1) champs \uline{localement}  constants ne dépendent que de

\margin{48}
